% Generated draft: requires mathematical and source review.
\documentclass[11pt]{article}
\usepackage[margin=25mm]{geometry}
\usepackage{fontspec}
\setmainfont{Latin Modern Roman}
\usepackage{amsmath,amssymb,mathtools,unicode-math}
\setmathfont{Latin Modern Math}
\usepackage{enumitem,needspace,xurl,hyperref}
\setlength{\parindent}{0pt}
\setlength{\parskip}{5pt}
\setlength{\emergencystretch}{3em}
\newcommand{\N}{\mathbb{N}}
\newcommand{\Z}{\mathbb{Z}}
\newcommand{\Q}{\mathbb{Q}}
\newcommand{\R}{\mathbb{R}}
\newcommand{\C}{\mathbb{C}}
\newcommand{\F}{\mathbb{F}}
\newcommand{\A}{\mathbb{A}}
\newcommand{\PP}{\mathbb{P}}
\newcommand{\E}{\mathbb{E}}
\newcommand{\eps}{\varepsilon}
\newcommand{\dd}{\,\mathrm{d}}
\begin{document}
\section*{UnsolvedMath ID: 6800006}
\textbf{Source problem code:} catalogue entry
\subsection{Definitions and mathematical statement}
In the usual flag-geometry framework, let \(G_{\C}\) be a complex semisimple Lie group and \(P\) a parabolic subgroup, so \(G_{\C}/P\) is a flag manifold. Let \(G_{\R}\) be a real form, with Lie algebra satisfying \(\mathfrak g_{\R}\otimes_{\R}\C\cong\mathfrak g_{\C}\). For a flag \(x\), the orbit
\[
\mathcal O=G_{\R}\cdot x\cong G_{\R}/\operatorname{Stab}_{G_{\R}}(x)
\]
is the proposed local model.

A manifold modeled on this orbit in the usual \((G_{\R},\mathcal O)\)-sense has an atlas of local diffeomorphisms to open subsets of \(\mathcal O\), with changes of coordinates given by restrictions of elements of \(G_{\R}\). Thus its dimension equals \(\dim_{\R}\mathcal O\). Equivalently, the standard developing/holonomy data consist of a local diffeomorphism
\[
\operatorname{dev}:\widetilde M\to\mathcal O
\]
and a homomorphism \(\rho:\pi_1(M)\to G_{\R}\) satisfying
\[
\operatorname{dev}(\gamma z)=\rho(\gamma)\operatorname{dev}(z).
\]
A global quotient of the entire orbit is a special case, not the definition of all such structures; completeness, injectivity of the developing map, and discreteness of holonomy are not imposed by the question.

Falbel asks which manifolds admit these modeled structures. A definite classification must fix, or explicitly range over, the complex group, flag type, real form, orbit, and manifold category. The extracted question does not choose compact versus noncompact manifolds or a classification up to diffeomorphism versus equivalence of geometric structures. The displayed framework makes the conventional local-model interpretation precise without claiming it is the only possible intended flag setting. The supplied research summary calls this a broad open structural direction and reports no complete classification located.
\subsection{Short English statement}
Classify manifolds admitting geometric charts modeled on a real-form orbit in a flag manifold, with transitions in that real form. The group, flag type, orbit, and global manifold restrictions must be specified.
\subsection{Sources}
\begin{itemize}
\item Local UnsolvedMath registry, permanent ID \(6800006\), exact source code \texttt{catalogue entry}, ``Manifolds modelled on flag manifolds.'' Proposer Elisha Falbel, Sorbonne Université, Paris; proposed year 2018.
\item Morgan and Pansu, \emph{A List of Open Problems in Differential Geometry} (2018), Question environment 6; \url{https://www.imo.universite-paris-saclay.fr/~pansu/problems_MTDG.tex}. Source-TeX extraction; recorded access 2026-07-29.
\item Supplied OPEN-TRIAGE research summary mentions Falbel--Guilloux--Will, spherical CR geometry, and complex/quaternionic settings without a complete citation or precise classification theorem.
\end{itemize}
\end{document}
